% © 2026 Ing. David Jaroš — CC BY-NC-ND 4.0
%
% This work is licensed under a Creative Commons Attribution-NonCommercial-NoDerivatives
% 4.0 International License (CC BY-NC-ND 4.0).
%
% See LICENSE.md for full license text.

\documentclass[12pt,a4paper]{article}
\usepackage[a4paper,margin=2.5cm]{geometry}
\usepackage{amsmath,amssymb,tcolorbox}
\tcbuselibrary{skins,breakable}
\newtcolorbox{statusbox}[1][]{colback=gray!8!white,colframe=black!60,arc=3pt,boxrule=1pt,title=Status Summary,#1}
\title{EW3.A: Killing-form / induced-inner-product analysis of $\theta_W$}
\author{Ing.\ David Jaroš}
\date{2026-05-12}
\begin{document}
\maketitle

Define induced product from kinetic term:
\[
\langle A,B\rangle_\Theta = \mathrm{Tr}[A,\Theta][B,\Theta]^\dagger.
\]
Naive fundamental-doublet normalisation gives equal norms for SU(2) and U(1):
\[
\langle Y^2\rangle \sim \langle \tau^2\rangle \Rightarrow \theta_W=\pi/4,
\]
which is phenomenologically excluded.

Using full fermion-content weighting (same counting as EW1):
\[
\frac{\langle Y^2\rangle}{\langle T_3^2\rangle}=\frac{5}{3}
\Rightarrow
\tan^2\theta_W=\frac{3}{5}
\Rightarrow
\theta_W=\arctan\sqrt{3/5}.
\]
Hence geometrical-angle interpretation is representation-content dependent.

\begin{statusbox}
\textbf{Hlavní výsledek}: Samotná naivní vnitřní metrika dává $\theta_W=\pi/4$ (NO-GO); realistický poměr vzniká až po započtení plného fermionického spektra.\\
\textbf{Klasifikace}: [MC/CONDITIONAL].\\
\textbf{Dopad na teorii}: Geometrická interpretace je možná, ale není čistě algebraická invarianta bez fermionového obsahu.\\
\textbf{Příští krok}: Explicitně odvodit váhové koeficienty reprezentací z UBT bez externích vstupů.\\
\textbf{Alpha je odvozena}: NE (unconditional)
\end{statusbox}

\end{document}
